%%%%%%%%%%%%%%%%%%%%%%%%%%%%%%%%%%%%%%%%%%%%%%%%%%%%%%%%%%%%%%%%%%%%%%%%%%%%%%
%  Partie 3 -- Marketing : promesses, coût face à l'abonnement, terrains d'usage
%  (3 diapositives)
%  Astuce de mise en page : les cartes à hauteur fixe se règlent avec carte[height=...,valign=top] ;
%  chaque icône est centrée dans un cadre invisible de 1,2 cm de haut pour aligner les titres.
%%%%%%%%%%%%%%%%%%%%%%%%%%%%%%%%%%%%%%%%%%%%%%%%%%%%%%%%%%%%%%%%%%%%%%%%%%%%%%
\section{Marketing}

%---------------------------------------------------------------------------
% Diapositive 10 : les trois promesses + slogans
%---------------------------------------------------------------------------
\begin{frame}{Trois promesses~: prix, autonomie, zones blanches}
  \begin{columns}[T]
    \begin{column}{0.325\textwidth}
      \begin{carte}[height=4.05cm,valign=top]
        \centering
        \begin{tikzpicture}\path (-0.7,-0.6) rectangle (0.7,0.6);\pic[insarouge,scale=1.7] at (0,0) {icone euro};\end{tikzpicture}\\[3pt]
        \gros{Petit prix}\\[3pt]
        {\footnotesize ≈ 69\,€ une seule fois. Aucun abonnement, aucune facture mensuelle, aucun frais caché.}
      \end{carte}
    \end{column}
    \begin{column}{0.325\textwidth}
      \begin{carte}[height=4.05cm,valign=top]
        \centering
        \begin{tikzpicture}\path (-0.7,-0.6) rectangle (0.7,0.6);\pic[insarouge,scale=1.9] at (0,0) {icone cadenas};\end{tikzpicture}\\[3pt]
        \gros{Autonomie}\\[3pt]
        {\footnotesize Ton antenne, ton logiciel libre, ta clé~: ni Orange, ni Free, ni SFR, ni réseau social.}
      \end{carte}
    \end{column}
    \begin{column}{0.325\textwidth}
      \begin{carte}[height=4.05cm,valign=top]
        \centering
        \begin{tikzpicture}\path (-0.7,-0.6) rectangle (0.7,0.6);\pic[scale=1.6] at (0,0) {icone mont};\end{tikzpicture}\\[3pt]
        \gros{Zones blanches}\\[3pt]
        {\footnotesize Campagne, montagne, zones de guerre, catastrophes naturelles~: là où le réseau manque ou tombe.}
      \end{carte}
    \end{column}
  \end{columns}
  \vspace{0.15cm}
  \begin{idee}\centering
    {\sffamily\bfseries\color{insarouge}Pas d'abonnement. Pas de compte. Pas de permission.}\\[3pt]
    {\sffamily\bfseries\color{insarouge}Là où le réseau s'arrête, le tien commence.}
  \end{idee}
\end{frame}

%---------------------------------------------------------------------------
% Diapositive 11 : coût cumulé face à l'abonnement
%---------------------------------------------------------------------------
\begin{frame}{Contre l'abonnement~: le coût sur trois ans}
  \begin{columns}[T]
    \begin{column}{0.62\textwidth}
      \centering
      \begin{tikzpicture}
        \begin{axis}[
            width=7.3cm,height=4.9cm,xmin=0,xmax=36,ymin=0,ymax=1180,
            xtick={0,12,24,36},xticklabels={départ,1 an,2 ans,3 ans},
            ytick={0,250,500,750,1000},
            tick label style={font=\scriptsize},
            ylabel={coût cumulé (€)},ylabel style={font=\scriptsize,yshift=-3pt},
            axis lines=left,grid=major,grid style={insagris!25,thin},
            legend style={at={(0.02,1.0)},anchor=north west,font=\scriptsize,draw=none,fill=none,
                          legend cell align=left},
            clip=false]
          \addplot[insagris,very thick] coordinates {(0,0) (36,643.68)};
          \addlegendentry{Forfait mobile moyen}
          \addplot[insaorange,very thick] coordinates {(0,363.3) (36,827.4)};
          \addlegendentry{Balise + abonnement}
          \addplot[insarouge,line width=2.2pt] coordinates {(0,69) (36,69)};
          \addlegendentry{Kit \nomproduit{}}
          \node[anchor=west,font=\scriptsize\bfseries,text=insagris] at (axis cs:36.6,644) {644\,€};
          \node[anchor=west,font=\scriptsize\bfseries,text=insaorange] at (axis cs:36.6,827) {827\,€};
          \node[anchor=west,font=\scriptsize\bfseries,text=insarouge] at (axis cs:36.6,69) {69\,€};
        \end{axis}
      \end{tikzpicture}
    \end{column}
    \begin{column}{0.35\textwidth}
      \begin{remarque}[Un complément]
        \begin{puces}
          \item Il ne remplace pas ton téléphone~: appels, internet.
          \item Aucun centre de secours n'est relié~: ce n'est pas une balise de détresse.
          \item Tes contacts doivent être équipés aussi.
        \end{puces}
      \end{remarque}
    \end{column}
  \end{columns}
  \vspace{0.1cm}
  \begin{pointcle}\centering\footnotesize
    69\,€, c'est moins de quatre mois d'un forfait mobile moyen (17,88\,€ par mois).
  \end{pointcle}
  \srcnote{Hypothèses~: 36 mois, 1\,\$ = 0,86\,€, hors remplacement de batterie. Balise~: inReach Mini 2 (328,90\,€), forfait 14,99\,\$ par mois, activation 39,99\,\$.
           Sources~: \cite{arcep_prix,garmin_mini2,garmin_plans}}
\end{frame}

%---------------------------------------------------------------------------
% Diapositive 12 : quatre terrains d'usage
%---------------------------------------------------------------------------
\begin{frame}{Là où le réseau s'arrête~: quatre terrains}
  \begin{columns}[T]
    \begin{column}{0.245\textwidth}
      \begin{carte}[height=4.0cm,valign=top]
        \centering
        \begin{tikzpicture}\path (-0.7,-0.6) rectangle (0.7,0.6);
          \pic[insagris,scale=1.05] at (-0.12,-0.05) {icone tour};\pic[scale=0.95] at (0.42,0.32) {icone croix};\end{tikzpicture}\\[3pt]
        \gros{Campagne}\\[3pt]
        {\footnotesize Il reste des zones blanches malgré 99\,\% de couverture 4G.}
      \end{carte}
    \end{column}
    \begin{column}{0.245\textwidth}
      \begin{carte}[height=4.0cm,valign=top]
        \centering
        \begin{tikzpicture}\path (-0.7,-0.6) rectangle (0.7,0.6);\pic[scale=1.5] at (0,0) {icone mont};\end{tikzpicture}\\[3pt]
        \gros{Montagne}\\[3pt]
        {\footnotesize 27 millions de Français pratiquent la randonnée~; le relief coupe le signal.}
      \end{carte}
    \end{column}
    \begin{column}{0.245\textwidth}
      \begin{carte}[height=4.0cm,valign=top]
        \centering
        \begin{tikzpicture}\path (-0.7,-0.6) rectangle (0.7,0.6);\pic[insagris,scale=1.5] at (0,0) {icone immeuble};\end{tikzpicture}\\[3pt]
        \gros{Guerre}\\[3pt]
        {\footnotesize Antennes détruites, coupures imposées~: un maillage n'a pas de centre.}
      \end{carte}
    \end{column}
    \begin{column}{0.245\textwidth}
      \begin{carte}[height=4.0cm,valign=top]
        \centering
        \begin{tikzpicture}\path (-0.7,-0.6) rectangle (0.7,0.6);\pic[insarouge,scale=1.8] at (0,0) {icone vague};\end{tikzpicture}\\[3pt]
        \gros{Catastrophes}\\[3pt]
        {\footnotesize Tempête, séisme, inondation~: les relais tombent, les kits sur batterie restent.}
      \end{carte}
    \end{column}
  \end{columns}
  \vspace{0.15cm}
  \begin{pointcle}\centering\footnotesize
    Un seul produit pour quatre terrains~: la même antenne, le même logiciel.
  \end{pointcle}
  \srcnote{Sources~: \cite{arcep_mobile,ffrando,laresilience}}
\end{frame}
